% command to make sure colorboxes don't create whitespace between code lines
% Reduce the space inside colorboxes to avoid extra whitespace between lines
\fboxsep=.6pt

% default settings for listings
\lstset{style=python}

\title[Daten Vorhersagen]{Daten Vorhersagen}

\subtitle{Lineare Regression}

\author[Wendl, Cyril] % (optional)
{Cyril Wendl}

\institute[KLW] % (optional)
{
    Kantonsschule im Lee
}

\date[\today] % (optional)

\logo{\includegraphics[width=2cm]{\thelogo}}

\titlegraphic{\vspace{-.7cm}\includegraphics[width=\textwidth]{"GrundlagenInfo/06_Datenbanken/Figures/title-emojis"}}

\frame{\titlepage}


\begin{frame}{Themenübersicht}
    \begin{enumerate}
        \item[\faCheckSquare{}] Daten verstehen (analysieren und visualisieren)
        \item[\faSquare] \textbf{Zahlen vorhersagen}
        \item[\faSquare] Kategorien vorhersagen
    \end{enumerate}
\end{frame}

\begin{frame}{Maschinelles Lernen}
    \begin{figure}[H]
        \only<1->{\centering\includegraphics[width=0.26\pagewidth]{"GrundlagenInfo/10_AusDatenLernen/Figures/intro_spam.jpg"}}
        \only<2->{\centering
\includegraphics[width=0.26\pagewidth]{"GrundlagenInfo/10_AusDatenLernen/Figures/intro_YT.jpg"}}
        \only<3->{\centering\includegraphics[width=0.26\pagewidth]{"GrundlagenInfo/10_AusDatenLernen/Figures/intro_scan.jpg"}}
        % \only<4->{\centering\includegraphics[width=0.3\pagewidth]{"GrundlagenInfo/10_AusDatenLernen/Figures/intro_recommendation.jpg"}}
    \end{figure}
\end{frame}

\begin{frame}{Zahlen vorhersagen mit Regression}{Unabhängige und abhängige Variablen}
    \begin{table}
        \centering
        \begin{tabular}{rcl}
            \only<1->{\textbf{Unabhängige Variable}} & \only<1->{\phantom{$\rightarrow$}} & \only<1->{\textbf{Abhängige Variable}} \\ \midrule\midrule
            \only<2->{Lernzeit}                      & \only<2->{$\rightarrow$}           & \only<2->{Note in der Prüfung}         \\ \midrule
            \only<3->{Alter eines Autos}             & \only<3->{$\rightarrow$}           & \only<3->{Wiederverkaufswert}          \\ \midrule
            \only<4->{Anzahl Schritte pro Tag}       & \only<4->{$\rightarrow$}           & \only<4->{Kalorienverbrauch}           \\ \midrule
            \only<5->{Aussentemperatur}              & \only<5->{$\rightarrow$}           & \only<5->{Eisverkaufszahlen}           \\ \midrule
            \only<6->{Anzahl Likes auf Social Media} & \only<6->{$\rightarrow$}           & \only<6->{Reichweite eines Posts}      \\
        \end{tabular}
    \end{table}
\end{frame}

\begin{frame}{\textbf{Unabhängige} und \textbf{abhängige Variablen}}
    \begin{table}
        \centering
        \begin{tabular}{p{3cm}|p{5cm}}
            \textbf{Unabhängige Variable} & Kontrollierbare Grösse, beeinflusst andere Variable(n).               \\ \midrule
            \textbf{Abhängige Variable}   & Messbare Grösse, wird beeinflusst von \textit{unabhängiger} Variable. \\
        \end{tabular}
    \end{table}
\end{frame}

\begin{frame}{Zahlen vorhersagen mit Regression}{Unabhängige und abhängige Variablen}
    \begin{figure}[H]
        \centering
        \includegraphics[width=0.5\pagewidth]{"GrundlagenInfo/10_AusDatenLernen/Figures/schlafdauer_vs_notenschnitt.pdf"}
    \end{figure}
    \begin{itemize}[<+->]
        \item Unabhängige Variable: Schlafdauer
        \item Abhängige Variable: Notenschnitt
        \item Ziel: Vorhersage des Notenschnitts anhand der Schlafdauer
        \item Beispiel: 7.3 Stunden Schlaf $\rightarrow$ Notenschnitt 5.2
    \end{itemize}
\end{frame}

\begin{frame}{Lineare Funktion}
    \begin{columns}
        \column{0.4\textwidth}
        Einfache Idee: zwei Datenpunkte verbinden, um eine Gerade zu berechnen:

        \begin{table}
            \centering
            \begin{tabular}{p{1.5cm}p{1.5cm}}
                \toprule
                \textbf{Schlaf-\newline Dauer\newline(=x)} & \textbf{Noten-\newline Schnitt\newline(=y)} \\
                \midrule
                3.54                                       & 3.57                                        \\
                9.617                                      & 6.0                                         \\
                \bottomrule
            \end{tabular}
        \end{table}

        \[m = \frac{y_2 - y_1}{x_2 - x_1}\]

        \[q = y_1 - m \cdot x_1\]

        \column{0.7\textwidth}
        \only<1>{
            \begin{figure}[H]
            \centering
            \includegraphics[width=\textwidth]{"GrundlagenInfo/10_AusDatenLernen/Figures/schlafdauer_vs_notenschnitt.pdf"}
        \end{figure}
        }
        \only<2>{
            \begin{figure}[H]
                \centering
                \includegraphics[width=\textwidth]{"GrundlagenInfo/10_AusDatenLernen/Figures/schlafdauer_vs_notenschnitt_2p_reg.pdf"}
            \end{figure}
        }
    \end{columns}
\end{frame}

\begin{frame}
    \begin{table}
        \footnotesize
        \centering
        \begin{tabular}{rrrrr}
            \toprule
            Punkt   & Schlafdauer & Notenschnitt & Vorhersage & Unterschied \\
            \midrule
            Punkt 1 & 5.39        & 4.61         & 4.61       & 0.00        \\
            Punkt 2 & 7.06        & 6.00         & 5.52       & 0.48        \\
            Punkt 3 & 7.87        & 5.95         & 5.95       & 0.00        \\
            \bottomrule
        \end{tabular}
    \end{table}

    \begin{figure}[H]
        \centering
        \includegraphics[width=0.6\pagewidth]{"GrundlagenInfo/10_AusDatenLernen/Figures/schlafdauer_vs_notenschnitt_3points.pdf"}
    \end{figure}
    Fehler:
    $$0^2+(-0.484)^2+0^2= 0.235$$
\end{frame}

\begin{frame}{Lineare Regression mit \textbf{OLS} - Ordinary Least Squares (Kleinste Quadrate)}
    $$m = \frac{n\sum x_iy_i - \sum x_i \sum y_i}{n\sum x_i^2 - (\sum x_i)^2}$$
    Das berechnete $m$ ist die Steigung der Geraden, die den Fehler minimiert. Diese kann nun in die Gleichung für $q$ eingesetzt werden:
    $$q = \frac{\sum y_i - m\sum x_i}{n}$$
\end{frame}

\begin{frame}
    \begin{figure}[H]
        \centering
        \includegraphics[width=.8\pagewidth]{"GrundlagenInfo/10_AusDatenLernen/Figures/schlafdauer_vs_notenschnitt_ols.pdf"}
    \end{figure}
\end{frame}

\begin{frame}{Multiple lineare Regression}{Mehrere unabhängige Variablen}
    \begin{columns}
        \column{0.5\textwidth}
        \begin{itemize}
            \item Beispiel: Note über Lernaufwand und Schlaf
            \item Erklärung:
                  \begin{itemize}
                      \item $x_1$: Lernaufwand
                      \item $x_2$: Schlaf
                      \item $y$: Note
                  \end{itemize}
        \end{itemize}
        \begin{block}{Formeln}
            \begin{itemize}
                \item Einfach: \\
                      $y = q + mx$
                \item Multipel: \\
                      $y = a_0 + a_1x_1 + a_2x_2$
            \end{itemize}
        \end{block}
        \column{0.65\textwidth}
        \begin{figure}[H]
            \centering
            \includegraphics[width=\textwidth]{"GrundlagenInfo/10_AusDatenLernen/Figures/linreg_3d.pdf"}
        \end{figure}
    \end{columns}
\end{frame}

\begin{frame}{Auftrag}{Jupyter Notebooks}
    \begin{itemize}
        \item Aufgabenblatt 3a durchlesen und vervollständigen
        \item Aufgabenblatt 3b lösen
    \end{itemize}
\end{frame}